%% kapa: exam=1 round=28 year=2017 period=1 form=A subject=real_estate subject_ko=부동산학원론 seq=7 number=87
%% source: Q-Net 감정평가사 기출문제 (cst00309) artlSeq=5199603 file=extracted/1008/5199603/2017년 제28회 감정평가사 1차시험 문제지/1교시 A형.pdf
%% conversion: kordoc 4.24.51 -> tools/gen1.py (auto); review=proofread-sonnet; flags=none
\begin{question}{87}
한국주택금융공사법에 의한 주택담보노후연금에 관한 설명으로 옳지 않은 것은?
\begin{choices}
  \item 단독주택, 다세대주택, 오피스텔, 상가주택 등이 연금의 대상주택이 된다.
  \item 연금 수령 중 담보 주택이 주택재개발, 주택재건축이 되더라도 계약을 유지할 수 있다.
  \item 연금의 방식에는 주택소유자가 선택하는 일정기간 동안 노후생활자금을 매월 지급 받는 방식이 있다.
  \item 가입자와 그 배우자는 종신거주, 종신지급이 보장되며, 가입자는 보증료를 납부해야 한다.
  \item 연금의 방식에는 주택소유자가 생존해 있는 동안 노후생활자금을 매월 지급 받는 방식이 있다.
\end{choices}
\end{question}
\answer{1}
